\begin{proof}
\textbf{1. Standing hypotheses and the degenerate case.}  
We assume \(X\ge 0\) a.s. and \(a>0\).  If \(\mathbb{P}(X=0)=1\) then
\(\mathbb{E}[X]=0\) and \(\mathbb{P}(X\ge a)=0\); thus (M) holds with equality.
Henceforth we suppose \(\mathbb{P}(X>0)>0\).

\textbf{2. Indicator of the event \(\{X\ge a\}\).}  
Define the indicator
\[
I(\omega):=\mathbf 1_{\{X\ge a\}}(\omega)=
\begin{cases}
1, & X(\omega)\ge a,\\[2pt]
0, & X(\omega)< a .
\end{cases}
\]
By definition of expectation of an indicator,
\begin{align}
\mathbb{E}[I]=\mathbb{P}(X\ge a).\tag{2.1}
\end{align}

\textbf{3. Pointwise bound.}  
For every \(\omega\in\Omega\),
\begin{align}
I(\omega)\le\frac{X(\omega)}{a}.\tag{3.1}
\end{align}
Indeed, if \(X(\omega)\ge a\) then \(I(\omega)=1\le X(\omega)/a\); if
\(X(\omega)<a\) then \(I(\omega)=0\le X(\omega)/a\).

\textbf{4. Taking expectations.}  
Since both sides of (3.1) are non‑negative random variables, monotonicity of
expectation yields
\begin{align}
\mathbb{E}[I]\le\mathbb{E}\!\left[\frac{X}{a}\right].\tag{4.1}
\end{align}
Linearity of expectation gives
\begin{align}
\mathbb{E}\!\left[\frac{X}{a}\right]=\frac{1}{a}\,\mathbb{E}[X].\tag{4.2}
\end{align}
Combining (4.1) and (4.2) we obtain
\begin{align}
\mathbb{E}[I]\le\frac{\mathbb{E}[X]}{a}.\tag{4.3}
\end{align}

\textbf{5. Replacing \(\mathbb{E}[I]\) by the probability.}  
Using (2.1) in (4.3) we get
\[
\mathbb{P}(X\ge a)\le\frac{\mathbb{E}[X]}{a},
\]
which is precisely Markov’s inequality (M).

\textbf{6. Equality case.}  
Equality in (M) can occur only if equality holds in every step above.
Equality in (4.1) requires equality in the pointwise bound (3.1) on a set of
full probability, because for non‑negative random variables \(Y\le Z\) a.s.
\[
\mathbb{E}[Y]=\mathbb{E}[Z]\;\Longrightarrow\;Y=Z\ \text{a.s.}
\tag{6.1}
\]
Applying (6.1) with \(Y=I\) and \(Z=X/a\) yields
\[
\frac{X}{a}=I\quad\text{a.s.}
\]
Hence, for \(\mathbb{P}\)-a.e. \(\omega\),
\[
X(\omega)=a\,I(\omega)=
\begin{cases}
a, & X(\omega)\ge a,\\[2pt]
0, & X(\omega)< a .
\end{cases}
\]
Thus \(X\) takes only the two values \(0\) and \(a\); equivalently,
\(X\) is a Bernoulli‑type random variable with
\[
\mathbb{P}(X=a)=\mathbb{P}(X\ge a),\qquad
\mathbb{P}(X=0)=1-\mathbb{P}(X\ge a).
\]
In this situation \(\mathbb{E}[X]=a\,\mathbb{P}(X\ge a)\), so the inequality
becomes an equality.

\textbf{7. Additional remarks.}  
If \(\mathbb{E}[X]=\infty\) the right‑hand side of (M) is infinite and the
inequality is trivial.  The requirement \(a>0\) is essential; for \(a\le0\)
the statement either reduces to \(\mathbb{P}(X\ge a)=1\) or becomes undefined.

Therefore Markov’s inequality holds for every non‑negative random variable,
with equality exactly when \(X\) assumes only the values \(0\) and \(a\) a.s.
\(\square\)
\end{proof}